% Source-derived chapter 08: Orthogonal motives of rank $2$
% Original source: standard-conjectures-abelian-fourfolds
\chapter{Orthogonal motives of rank $2$}
\label{standard-conjectures-abelian-fourfolds-chapter-08}
\source{standard-conjectures-abelian-fourfolds}

\begin{remark}[Source section]
\label{standard-conjectures-abelian-fourfolds-chapter-08-source}
\source{standard-conjectures-abelian-fourfolds}
\source{standard-conjectures-abelian-fourfolds}
This chapter reproduces the corresponding section of the retrieved
LaTeX source, including its definitions, statements, examples, and
proofs. It has not yet been translated into Lean.
\notready
\end{remark}

% The source section title was: Orthogonal motives of rank $2$
We start by stating our main technical result,  Theorem \ref{mainthm}. The remark and proposition right after will hopefully give some intuitions on the hypothesis of the statement. We conclude the section by showing that Theorem \ref{mainthm} implies  Theorem \ref{thmscht}. The proof of Theorem \ref{mainthm} will take the rest of the paper and will be ended in Section 13.
\begin{theorem}
\source{standard-conjectures-abelian-fourfolds}
\notready\label{mainthm}
\source{standard-conjectures-abelian-fourfolds}
Let $K$ be a $p$-adic field, $W$  its ring of integers and $k$  its residue field. Let us fix an embedding $\sigma : K \hookrightarrow \C$.
Let \[M \in \CHM(W)_\Q\] be a motive in mixed characteristic and consider the motives induced by pullbacks $M_{\vert\C} \in \CHM(\C)_\Q$ and $M_{\vert{k}} \in \CHM(k)_\Q$. Consider $V_B$ and $V_Z$ the two $\Q$-vector spaces defined as
\[V_B = R_B(M_{\vert\C} ) \hspace{0.5cm}
%(the Betti realization of $M_{\vert\C}$) 
\textrm{and} \hspace{0.5cm}
V_Z = \Hom_{\NUM(k)_\Q}(\one, M_{\vert{k}}).\]
%(with $\NUM(k)_\Q$ being the category of numerical motives over $k$).

Let $q$ be a quadratic form on $M$, by which we mean a morphism in $ \CHM(W)_\Q$  of the form
$q: \Sym^2 M  \longrightarrow \one.$
 Consider the two $\Q$-quadratic forms induced by $q$ on $V_B$ and $V_Z$ respectively, 
\[q_B=R_B(q)  \hspace{0.5cm} \textrm{and}  \hspace{0.5cm} q_Z(\cdot)= (q_{\vert{k}} \circ \Sym^2(\cdot)).\]
Suppose that the following holds:
\begin{enumerate}

\item\label{assumptiondim2} The two $\Q$-vector spaces $V_B$  and $V_Z$ are of dimension $2$.
\source{standard-conjectures-abelian-fourfolds}

\item The pairing $q_B$ on $V_B$ is a polarization of Hodge structures.

\end{enumerate}


Then the quadratic form $q_Z$ is positive definite.

\end{theorem}

\begin{remark}
\source{standard-conjectures-abelian-fourfolds}
\notready\label{hypothesismainthm}
\source{standard-conjectures-abelian-fourfolds}
\begin{enumerate}
\item The main example of such a motive $M$ we have in mind is an exotic motive of an abelian fourfold (see Section \ref{sectionexotic}). Another example coming from geometry is given in Proposition \ref{propfermat}.
\item\label{homvsrat} One can actually make the  hypothesis a little bit more flexible and work with homological motives instead of Chow motives. This will not matter for our application to abelian varieties.
\source{standard-conjectures-abelian-fourfolds}
\end{enumerate}
\end{remark}
\begin{prop}
\source{standard-conjectures-abelian-fourfolds}
\notready\label{tortura}
\source{standard-conjectures-abelian-fourfolds}
Under the hypothesis of Theorem \ref{mainthm} the following holds:
\begin{enumerate}
\item The quadratic form $q_Z$   on $V_Z$ is non-degenerate. 
\item  Given a classical realization $R$, the vector space $R(M_{\vert{k}})$  is of dimension two.
\item Given a classical realization $R$, the vector space $R(M_{\vert{k}})$ is spanned by algebraic classes.
\item Numerical equivalence on $\Hom_{\CHM(k)_\Q}(\one,M_{\vert{k}})$ coincides with the homological equivalence for any classical cohomology.
\item\label{giorno} Fix a classical realization $R$ and call $L$ the field of coefficients of  $R$. Then the equality $V_Z\otimes _\Q L = R(M)$ holds.
\source{standard-conjectures-abelian-fourfolds}

\end{enumerate}
\end{prop}
\begin{proof}
Let us start with (2). By the comparison theorem between singular and $\ell$-adic cohomology, we have ${\dim_\Q V_B = \dim_{\Q_\ell}R_\ell(M_{\vert{\C}})}$, for all primes $\ell$, including $\ell=p$.  Then, by smooth proper base change we have   $\dim_{\Q_\ell}R_\ell(M_{\vert{\C}})=\dim_{\Q_\ell}R_\ell(M_{\vert{k}})$ for all $\ell\neq p$ and finally by the $p$-adic comparison theorem we have
 $\dim_{\Q_p}R_p(M_{\vert{\C}})=\dim_{\Frac(W(k))}R_\crys(M_{\vert{k}})$. This concludes (2) as, by hypothesis, $\dim_\Q V_B=2$. 
 
 The proof of (3)-(5) goes as in Proposition \ref{motivisupersingolari}. There, the hypothesis that the motive was of abelian type\footnote{For our main application the motive $M$ will be anyway of abelian type.} was used to ensure that all  realizations have the same dimension (see Remark \ref{remarkrank}). Here, this is replaced by part (2).
 
 Let us now show part (1). First notice that it is enough to show that $R(q_{\vert{k}})$ is non-degenerate (for some classical realization) because of parts (2)-(5).
 Then, again by using the comparison theorems, this is equivalent to the fact that $R_B(q)$ is non-degenerate. On the other hand this is the case as $R_B(q)$ is a polarization of Hodge structures.
%Let us now assume (a), which means that $\dim_FR(M_{\vert{k}})= \dim V_Z$, and let us call $n$ this dimension.  Fix $n$ elements $v_1,\ldots,v_n$ in $R(M_{\vert{k}})$ so that their reduction modulo numerical equivalence is a $\Q$-basis of $V_Z$. These elements form an $F$-basis of $R(M_{\vert{k}})$. Indeed, by Lemma \ref{lemmanumerical} we can find algebraic classes $w_1,\ldots,w_n$ in $R(M_{\vert{k}})$ such that $R(q)(v_i,w_j)=\delta_{ij}$, which implies that there is no linear relation between $v_1,\ldots,v_n$. This implies (b), which clearly implies (c).
%Let us now assume (c) and fix $v_1,\ldots,v_n$ a set of algebraic classes which form a basis of $R(M_{\vert{k}})$. Let us call $V$ their $\Q$-span. Fix an algebraic class $w \in R(M_{\vert{k}})$. The pairing $R(q)(v_i,w)\in \Q \subset F,$ hence $w$ define an element of $V^*$ which by Lemma \ref{lemmanumerical} implies that $w \in V$. This means also that $R(M_{\vert{k}}) = V \otimes_\Q F$, which implies that homological and numerical equivalence coincide on $V$ as well.
\end{proof}
\begin{proof}[Proof of Theorem \ref{thmscht}]
Let $A$ be an abelian fourfold and $L$ be a hyperplane section. By Proposition \ref{reducefinite} we can suppose that $A$ and $L$ are defined over a finite field $k=\mathbb{F}_q$. Notice that, in order to prove Theorem \ref{thmscht}, we can (and will) replace $k$ by a finite extension.

Consider now the decomposition from Proposition \ref{primadecomposizione}(\ref{decomposizioneQ}) and, among the factors of this decomposition, consider 
those that are exotic (Definition \ref{defexotic}). Then there exist a finite extension $k'$ of $k$ and a  CM-structure for $A\times_k k'$ such that any exotic motive  of $\mathfrak{h}^{4}(A\times_k k')$ has dimension two (Proposition \ref{exoticdim2}).  By Theorem \ref{HondaTate}, the CM-structure can be lifted    to characteristic zero.  Moreover, by Proposition \ref{anypolarization}, it is enough to work with a single $L$ for a given abelian fourfold. By Corollary \ref{CMstab} we can choose such an $L$ so that it lifts to characteristic zero and the CM-structure is $*$-stable (Definition \ref{definitionCM}).

 
Now, by Lemma \ref{accoppiamenti1} we can work with the pairing $\langle \cdot, \cdot \rangle_{1,\mot}^{\otimes 4}$ instead of $\langle \cdot, \cdot \rangle_{4,\mot}$. For this pairing, the decomposition from Proposition \ref{primadecomposizione}(\ref{decomposizioneQ}) is orthogonal, hence we can work with a single motive of the decomposition. By Propositions \ref{oklefschetz} and \ref{talvez}(\ref{generatolefschetz}) we are reduced to motives that are exotic.

Finally, those exotic motives $\mathcal{M}_I$ are settled by Theorem \ref{mainthm} by setting $M=\mathcal{M}_I(2)$. Notice that all the hypothesis of Theorem \ref{mainthm} are satisfied. Indeed, the motive lives in mixed characteristic (Corollary \ref{lifting motives}), together with its quadratic form (because $L$ lifts to characteristic zero). The space $V_B$ is clearly of dimension two; so it is $V_Z$ by   Proposition  \ref{talvez}(\ref{generatoalgebrico}).  Last, the quadratic form $q_B=R_B(\langle \cdot, \cdot \rangle_{1,\mot}^{\otimes 4})$ is a polarization as $R_B(\langle \cdot, \cdot \rangle_{1,\mot})$ is so.
\end{proof}
